\documentclass[conference]{IEEEtran}
\IEEEoverridecommandlockouts
\usepackage[utf8]{inputenc}
\usepackage{pgfplots}
\DeclareUnicodeCharacter{2212}{−}
\usepgfplotslibrary{groupplots,dateplot}
\usetikzlibrary{patterns,shapes.arrows}
\pgfplotsset{compat=newest}
\usepackage{listings}
\usepackage{lipsum}
\usepackage{todonotes}
\usepackage{float} 
\usepackage{caption}
\usepackage{subcaption}
\usepackage{graphicx}
\usepackage{multirow}
\usepackage{hyperref}
\usepackage{booktabs}
%\usepackage{url}
\usepackage{tabularx}
\usepackage{xcolor}
\usepackage{soul}
\newcommand{\mathcolorbox}[2]{\colorbox{#1}{$\displaystyle #2$}}
\usepackage{todonotes}
% Unidades del SI
\usepackage{siunitx} 
\usepackage{tikz}
\usepackage{pgfplots}
\usepackage{pgfplotstable}
\pgfplotsset{compat=1.18}

\usepackage[acronym]{glossaries}
\makeglossaries

\newacronym{CNN}{CNN}{Convolutional Neural Network}
\newacronym{LSTM}{LSTM}{Long Short Term Memory}
\newacronym{NB}{NB}{Naive Bayes}
\newacronym{LLM}{LLM}{Large Language Models}
\newacronym{ECE}{ECE}{Expected Calibration Error}
\newacronym{TS}{TS}{Temperature Scaling}
\newacronym{FEINA}{FEINA}{Financial Education Corpus IN SpAnish}
\newacronym{AI}{AI}{Artificial Intelligence}
\newacronym{BERT}{BERT}{Bidirectional Encoder Representations from Transformers}

\newacronym{DNN}{DNN}{Deep Neural Network}

\newacronym{UQ}{UQ}{Uncertainty Quantification}

\newacronym{FD}{FD}{Feature Densities}
\newacronym{BETO}{BETO}{Spanish BERT}

\newacronym{NLP}{NLP}{Natural Language Processing}
\newacronym{JSD}{JSD}{Jensen-Shannon Distance}
\newacronym{UCCA}{UCCA}{Universal Cognitive Conceptual
Annotation}
\newacronym{CEM}{CEM}{Confidence Estimation Module}
\newacronym{WER}{WER}{Word Error Rate}
\newacronym{EASSE}{EASSE}{Easier Automatic Sentence Simplification Evaluation}

\newacronym{ASR}{ASR}{Automatic Speech Recognition}

\newacronym{SAMSA}{SAMSA}{Simplification Automatic evaluation Measure through Semantic Annotation}

\newacronym{MCD}{MCD}{Monte Carlo Dropout}
\newacronym{UE}{UE}{Uncertainty Estimation}
\newacronym{SARI}{SARI}{System output Against References
and against the Input sentence}
\newacronym{DEB}{DEB}{Deep Ensemble Based}
\newacronym{BNN}{BNN}{Bayesian Neural Networks}

\newacronym{BLEU}{BLEU}{Bilingual Evaluation Understudy Score}

\newacronym{GPT}{GPT}{Generative Pre-trained Transformer}

\newacronym{RAG}{RAG}{Retrieval-Augmented Generation}
\newacronym{QA}{QA}{Question Answering}


\ifCLASSOPTIONcompsoc
\usepackage[nocompress]{cite}
\else
\usepackage{cite}
\fi
\ifCLASSINFOpdf
\else
\fi
\hyphenation{op-tical net-works semi-conduc-tor}
\makeatletter
\newcommand{\linebreakand}{
  \end{@IEEEauthorhalign}
  \hfill\mbox{}\par
  \mbox{}\hfill\begin{@IEEEauthorhalign}
}
\providecommand{\keywords}[1]
{
  \small	
  \textbf{\textit{Keywords---}} #1
}
\makeatother

\renewcommand{\thefigure}{\arabic{figure}}

\begin{document}

% I need to upgrade this title
\title{A Benchmark of Large Language Models for Question and Answering with RAG applied using Syllabus Documents in Spanish}


% author names and affiliations
% use a multiple column layout for up to three different
% affiliations

\author{

\IEEEauthorblockN{Caleb Alfaro-Moreira}
\IEEEauthorblockA{\textit{Instituto Tecnológico de Costa Rica }\\
Alajuela, Costa Rica\\
calalfaro@estudiantec.cr}

\and

\IEEEauthorblockN{ Saúl Calderón-Ramírez}
\IEEEauthorblockA{\textit{Instituto Tecnológico de Costa Rica }\\
Cartago, Costa Rica\\
sacalderon@itcr.ac.cr}

\and

\IEEEauthorblockN{Martín Solís}
\IEEEauthorblockA{\textit{Instituto Tecnológico de Costa Rica }\\
Cartago, Costa Rica\\
marsolis@itcr.ac.cr}

\linebreakand

\IEEEauthorblockN{Mauricio Araya}
\IEEEauthorblockA{\textit{Instituto Tecnológico de Costa Rica }\\
Cartago, Costa Rica\\
marroyo@itcr.ac.cr}

\and

\IEEEauthorblockN{Mónica Hernández}
\IEEEauthorblockA{\textit{Instituto Tecnológico de Costa Rica (??)}\\
Cartago, Costa Rica\\
mohernandez@itcr.ac.cr}

\and

\IEEEauthorblockN{José Aguilar}
\IEEEauthorblockA{\textit{Instituto Tecnológico de Costa Rica (??)}\\
Cartago, Costa Rica\\
???}

}


\maketitle

\begin{abstract}

%TODO: write last, once results are in. Should cover: motivation
%(on-premise LLMs for university digital transformation), method (RAG +
%QA benchmark over course-program documents), models compared, and the
%headline result (which model/size offers the best accuracy-efficiency
%tradeoff).
    
\end{abstract}

\keywords{ \footnote{\textbf{}} }

\section{Introduction}

%TODO: expand into full introduction. Suggested flow:
%  1. Universities are adopting AI tools as part of digital transformation
%     strategies (admissions, student support, academic advising, etc.)
%  2. Most capable LLMs are proprietary and cloud-hosted, which raises
%     cost, data privacy, and infrastructural dependency concerns for
%     institutions with limited resources.
%  3. Smaller open-source LLMs are increasingly viable for self-hosted
%     deployment, but their reliability on domain-specific QA tasks
%     grounded in institutional documents is not well characterized.
%  4. This paper benchmarks three such models over a QA dataset derived
%     from university postgraduate course programs, using a RAG pipeline,
%     to assess whether they are accurate and efficient enough for
%     practical institutional deployment.
%  5. Brief roadmap of the paper's sections.

\section{Problem Definition}

Universities are increasingly incorporating \gls{AI} tools into administrative and academic processes as part of broader digital transformation strategies. However, the most capable \glspl{LLM} are typically proprietary and cloud-hosted, which raises concerns around recurring costs, data privacy, and dependency on external infrastructure --- concerns that are particularly relevant for institutions with limited computational resources. This creates a need to evaluate whether smaller, open-source \glspl{LLM}, deployable on modest institutional hardware, can perform reliably on domain-specific \gls{QA} tasks grounded in an institution's own documents, such as course syllabi and program descriptions.

This study benchmarks six open-source \gls{LLM} configurations, spanning two model families (Qwen2.5 and Llama 3.1 at the efficient tier, DeepSeek-R1-Distill across the mid and upper-bound tiers) at multiple parameter scales and quantization levels, using a \gls{RAG} pipeline over a curated set of university postgraduate course documents. Model performance is assessed across three question types (binary, short-answer, and open-ended), with the goal of identifying the configurations that offer the best practical balance between answer accuracy and computational efficiency for self-hosted deployment at the university.

%TODO: once results are available, close this section with explicit
%research questions, e.g.:
%  RQ1: How does QA accuracy vary across model and question type?
%  RQ2: What is the accuracy-vs-inference-cost tradeoff across models?
%  RQ3: Is any of the benchmarked models viable for deployment on the
%       university's available hardware?

\section{Experimental Design}

\subsection{Dataset}
The dataset was generated from 15 university postgraduate course program files in PDF and DOCX formats, all in Spanish. These documents were processed through a text-extraction and \gls{RAG} pipeline to create a retrieval context for the \glspl{LLM}.
A total of 270 questions were constructed, divided into three categories: 90 binary (yes/no) questions, 90 short-answer questions, and 90 open-ended questions.
The questions were generated using Gemini Notebook after an initial manual review of the source documents to identify key information and topics.
Each question was validated to ensure it could be answered based on the information contained within the source documents. 

\begin{itemize}
    \item 90 yes or no questions -> Metrics: F1 score
    \item 90 short answer questions -> Metrics: F1 score
    \item 90 open ended questions -> Metrics: BERT Score, ROUGE, METEOR
\end{itemize}

%TODO: describe the source documents (15 university postgraduate course
%program files, PDF/DOCX, in Spanish) and the text-extraction/RAG
%pipeline used to build the retrieval context, plus how the QA dataset
%(90 questions per category) was constructed/validated.

%TODO (Limitations): retrieval uses free (unfiltered) semantic search
%across all 15 documents, deliberately, to evaluate the full RAG
%pipeline rather than isolating generation quality. Manual inspection
%during pipeline validation found cases where the correct source chunk
%was retrieved but ranked below chunks from a different, topically
%adjacent document (e.g. a course's own curriculum sentence on SOC/AI
%outranked by a generic policy-governance document's incident-response
%language) -- consistent with known limitations of lightweight
%multilingual embedding models (paraphrase-multilingual-MiniLM-L12-v2)
%on nuanced domain content. This affects all benchmarked models equally,
%preserving fairness of the comparison, but likely places an upper bound
%on achievable accuracy independent of generator quality.

\subsection{LLMs to benchmark}

Models are grouped into three tiers, reflecting the accuracy-vs-efficiency tradeoff central to this study's motivation. All models are loaded from their native (safetensors) Hugging Face repositories and quantized on-the-fly via \texttt{bitsandbytes}, rather than relying on pre-quantized GGUF files, so that a single loading and generation codepath is shared across every configuration.

\begin{itemize}
  \item \textbf{Efficient tier} (target deployment candidates):
  \begin{itemize}
    \item Qwen2.5-3B-Instruct (full precision)
    \item Llama-3.1-8B-Instruct (4-bit)
  \end{itemize}
  \item \textbf{Mid tier}:
  \begin{itemize}
    \item Llama-3.1-8B-Instruct (full precision)
    \item DeepSeek-R1-Distill-Qwen-14B (4-bit)
    \item DeepSeek-R1-Distill-Qwen-14B (8-bit)
  \end{itemize}
  \item \textbf{Upper-bound reference tier} (to quantify what is given up by deploying smaller models):
  \begin{itemize}
    \item DeepSeek-R1-Distill-Qwen-14B (full precision)
  \end{itemize}
\end{itemize}

%TODO: A Llama-3.3-70B-Instruct configuration was initially considered as
%an additional upper-bound reference, but was excluded: on Kabre's
%single-GPU-per-node allocation (NVIDIA L40S, 46GB VRAM), a 4-bit
%quantization of a 70B model leaves too little headroom for the KV cache
%at the context lengths used here, and true multi-node distributed
%inference was judged out of scope given the project timeline.

%TODO: Qwen2.5-32B-Instruct (both 4-bit and 8-bit) was also considered,
%as a mid/upper-bound-tier counterpart to the DeepSeek-R1-Distill-14B
%configurations, but was excluded: downloading the ~65GB fp16 weights
%failed with a Lustre disk-quota error on Kabre's /data allocation
%(a per-user quota, not a lack of raw storage capacity on the mount).
%Requesting a quota increase from CeNAT was judged not worth the delay
%given the project timeline. As a result this study's mid/upper tiers
%are represented only by the DeepSeek-R1-Distill-14B family, and the
%comparison lacks a second model family at the 30B+ parameter scale.

For each configuration, model size (parameters), GPU memory footprint, inference latency, and approximate GPU cost in the current market are reported and compared. Cost is estimated as mean generation time per query multiplied by an hourly GPU rental rate, projected to a cost-per-1000-queries figure; this is not what Kabr\'e itself charges (allocation-based, not billed per hour) but a stand-in for what the same workload would cost on commercial cloud infrastructure --- the actual comparison point for this paper's cost argument against paid, proprietary LLM APIs. The rate used is \$0.96/hr, the on-demand dedicated price for an NVIDIA L40S --- the same GPU model used throughout this study's own infrastructure --- as listed by Spheron Network~\cite{spheron2026l40s}.

\subsection{Infrastructure}

Experiments are run on Kabr\'e, the national HPC cluster operated by Costa Rica's Centro Nacional de Alta Tecnolog\'ia (CeNAT), accessible to CONARE-affiliated universities including the Instituto Tecnol\'ogico de Costa Rica. Jobs are scheduled via Slurm on the \texttt{nukwa} partition family, which provides single-GPU compute nodes: 4 nodes with an NVIDIA Tesla V100 (32GB), and 7 nodes with an NVIDIA Ada L40S (46GB usable VRAM).

%TODO: V100 (Volta, compute capability 7.0) nodes are excluded from this
%study: the PyTorch build required for the other models' CUDA/driver
%stack (cu128) is compiled for compute capability >=7.5 and does not
%include kernels for Volta, causing a
%"no kernel image is available for execution on the device" error at
%runtime. All reported results therefore use L40S nodes exclusively.

This setup is itself part of the paper's motivation: rather than benchmarking on ad hoc consumer hardware, all models are evaluated on infrastructure a Costa Rican university can realistically access today through CeNAT, without relying on paid commercial cloud APIs.

\section{Results}

Tables and figures in this section are typeset directly from the result CSV files in \texttt{data/} at compile time via \texttt{pgfplotstable} and \texttt{pgfplots}, rather than transcribed by hand, so they always reflect the current contents of those files. \texttt{comparison\_accuracy.csv} holds all three question categories in one file (long format); the \texttt{\textbackslash FilterByCategory} macro below filters it per category at compile time, so no separate per-category files are needed.

\pgfplotstableread[col sep=comma]{data/comparison_accuracy.csv}\accuracytable
\pgfplotstableread[col sep=comma]{data/comparison_efficiency.csv}\efficiencytable

\makeatletter
% Loops over \accuracytable, building one tabular row (formatted with
% \pgfmathprintnumber) for every row whose 'category' column equals #1,
% using only the columns listed in #2 (comma-separated, first output
% column is always model_key). Output is accumulated into a macro and
% printed only after the loop's own grouping has closed — emitting
% "&"/"\\" directly from inside \foreach's body breaks tabular's
% alignment parser, hence the two-step accumulate-then-print structure.
\newcommand{\FilterByCategory}[2]{%
  \gdef\FBC@accum{}%
  \pgfplotstablegetrowsof{\accuracytable}%
  \pgfmathtruncatemacro{\FBC@numrows}{\pgfplotsretval-1}%
  \foreach \FBC@i in {0,...,\FBC@numrows}{%
    \pgfplotstablegetelem{\FBC@i}{category}\of\accuracytable%
    \edef\FBC@cat{\pgfplotsretval}%
    \def\FBC@target{#1}%
    \ifx\FBC@cat\FBC@target
      \pgfplotstablegetelem{\FBC@i}{model_key}\of\accuracytable%
      \edef\FBC@row{\pgfplotsretval}%
      \foreach \FBC@col in {#2}{%
        \pgfplotstablegetelem{\FBC@i}{\FBC@col}\of\accuracytable%
        \edef\FBC@val{\pgfplotsretval}%
        \xdef\FBC@row{\FBC@row\noexpand&\noexpand\pgfmathprintnumber[fixed,fixed zerofill,precision=3]{\FBC@val}}%
      }%
      \xdef\FBC@accum{\FBC@accum\FBC@row\noexpand\\}%
    \fi
  }%
  \FBC@accum
}
\makeatother

\subsection{Binary Questions}

\begin{table}[H]
\centering
\caption{Binary question F1 score, RAG vs. full-context, by model.}
\label{tab:binary-accuracy}
\begin{tabular}{lcc}
\toprule
Model & F1 (RAG) & F1 (Full-context) \\
\midrule
\FilterByCategory{binary}{f1_rag,f1_fullcontext}
\bottomrule
\end{tabular}
\end{table}

\subsection{Short-Answer Questions}

\begin{table}[H]
\centering
\caption{Short-answer question token-overlap F1 score, RAG vs. full-context, by model.}
\label{tab:short-answer-accuracy}
\begin{tabular}{lcc}
\toprule
Model & F1 (RAG) & F1 (Full-context) \\
\midrule
\FilterByCategory{short_answer}{f1_rag,f1_fullcontext}
\bottomrule
\end{tabular}
\end{table}

\subsection{Open-Ended Questions}

\begin{table}[H]
\centering
\caption{Open-ended question metrics (BERTScore F1, ROUGE-L, METEOR), RAG vs. full-context, by model.}
\label{tab:open-ended-accuracy}
\resizebox{\columnwidth}{!}{%
\begin{tabular}{lcccccc}
\toprule
Model & BERTScore (RAG) & BERTScore (Full-ctx) & ROUGE-L (RAG) & ROUGE-L (Full-ctx) & METEOR (RAG) & METEOR (Full-ctx) \\
\midrule
\FilterByCategory{open_ended}{bertscore_f1_rag,bertscore_f1_fullcontext,rouge_l_rag,rouge_l_fullcontext,meteor_rag,meteor_fullcontext}
\bottomrule
\end{tabular}
}
\end{table}

\subsection{Efficiency and Cost}

\begin{table}[H]
\centering
\caption{Mean generation time and estimated cost per 1000 queries, RAG vs. full-context, by model. Cost estimated at \$0.96/hr (NVIDIA L40S, on-demand)~\cite{spheron2026l40s}.}
\label{tab:efficiency}
\resizebox{\columnwidth}{!}{%
\pgfplotstabletypeset[
  columns={model_key,mean_generation_time_s_rag,mean_generation_time_s_fullcontext,est_cost_per_1000_queries_usd_rag,est_cost_per_1000_queries_usd_fullcontext},
  columns/model_key/.style={string type, column name={Model}},
  columns/mean_generation_time_s_rag/.style={column name={Time RAG (s)}, fixed, fixed zerofill, precision=3},
  columns/mean_generation_time_s_fullcontext/.style={column name={Time Full-ctx (s)}, fixed, fixed zerofill, precision=3},
  columns/est_cost_per_1000_queries_usd_rag/.style={column name={Cost RAG (\$/1k)}, fixed, fixed zerofill, precision=4},
  columns/est_cost_per_1000_queries_usd_fullcontext/.style={column name={Cost Full-ctx (\$/1k)}, fixed, fixed zerofill, precision=4},
  every head row/.style={before row=\toprule, after row=\midrule},
  every last row/.style={after row=\bottomrule},
]\efficiencytable
}
\end{table}

\begin{figure}[H]
\centering
\begin{tikzpicture}
\begin{axis}[
  ybar,
  bar width=6pt,
  width=\columnwidth,
  height=5cm,
  ylabel={Estimated cost per 1000 queries (USD)},
  symbolic x coords={deepseek-r1-14b-4bit,deepseek-r1-14b-8bit,deepseek-r1-14b-f16,llama3.1-8b-4bit,llama3.1-8b-f16,qwen2.5-3b},
  xtick=data,
  x tick label style={rotate=45, anchor=east, font=\tiny},
  legend style={font=\tiny, at={(0.5,1.15)}, anchor=north, legend columns=2},
  ymajorgrids=true,
]
\addplot table[x=model_key, y=est_cost_per_1000_queries_usd_rag, col sep=comma] {data/comparison_efficiency.csv};
\addplot table[x=model_key, y=est_cost_per_1000_queries_usd_fullcontext, col sep=comma] {data/comparison_efficiency.csv};
\legend{RAG, Full-context}
\end{axis}
\end{tikzpicture}
\caption{Cost per 1000 queries, RAG vs. full-context, by model. RAG is consistently cheaper across all six models, since it sends far fewer tokens per prompt than stuffing the entire 15-document corpus into context.}
\label{fig:cost-comparison}
\end{figure}

%TODO: add a second figure directly addressing the paper's central
%question — accuracy vs. cost tradeoff across models (e.g. binary F1 on
%the y-axis, cost per 1000 queries on the x-axis, one point per model,
%RAG and full-context as separate series) — once the specific framing
%for the Discussion section's argument is finalized.

\section{Discussion}

%TODO: synthesize the tables above. Candidate points, drawn from
%observations made during pipeline development and validation:
%  - RAG vs. full-context: does full-context's higher cost translate to
%    a proportionally better accuracy? (Table \ref{tab:binary-accuracy}
%    and \ref{tab:efficiency} suggest this is NOT uniform across models
%    or categories -- worth checking short-answer F1 in particular,
%    where full-context accuracy looks lower for several models despite
%    higher cost.)
%  - DeepSeek-R1-Distill's reasoning overhead: skip_reasoning=True was
%    used throughout (see model_config.py) after smoke testing showed
%    negligible accuracy impact but a ~5-10x latency cost when reasoning
%    was left enabled -- worth stating this explicitly as a design
%    decision with its own accuracy/cost tradeoff data if available.
%  - Retrieval quality ceiling: RAG's retrieval step was observed to
%    occasionally rank a topically-adjacent but incorrect document above
%    the correct source chunk (see Limitations note in Experimental
%    Design), which may partly explain any cases where full-context
%    outperforms RAG on accuracy despite retrieving from a much smaller,
%    cheaper context.
%  - Efficient-tier models (Qwen2.5-3B, Llama-3.1-8B) vs. larger
%    DeepSeek-R1-Distill-14B variants: does the smaller/cheaper tier
%    retain most of the larger tier's accuracy, supporting the paper's
%    central deployment-feasibility argument?

\section{Conclusions}

%TODO: write after results/discussion sections exist.

% Referencias
\bibliographystyle{ieeetr}
\bibliography{bibli}

\end{document}